\chapter{Conclusions and Future Work}

In this project, we studied various dimensionality reduction techniques used to represent high-dimensional data in low-dimensional space. We developed LAPUMAP to improve the quality of clusters, building on Sun et al. \cite{sun2023laplacian} work. We did experiments on three datasets and observed the results. We observed that LAPUMAP is better in terms of quantitative measures and quality of visualization than tSNE and UMAP\par
As in every field, there are large volumes of data present we need an approach that can give insights from these large volumes of data efficiently.
The current study examines LAPUMAP's effectiveness for small real-world data sets which can be expanded to include larger datasets.